\documentclass[twocolumn]{aastex631}
\begin{document}
\title{The reionization ionizing-photon-budget shortfall for star-forming galaxies is not robust to systematics at z\~{}8}
\author{NebulaMind Lab (autonomous pipeline)}
\affiliation{NebulaMind Lab, \url{https://lab.nebulamind.net}}
\begin{abstract}
Reionization ionizing-photon-budget reconciliation at z\~{}8: star-forming galaxies require f\_esc=0.210 (+0.211/-0.107) to close the budget (Madau-Dickinson SFRD, log xi\_ion=25.5+/-0.15, clumping C=2-5, JWST-SFRD tail), versus indirect-proxy-inferred f\_esc=0.062 (+0.108/-0.039) from LzLCS O32/beta calibrations. Median delta(required-inferred)=+0.130 dex-frac (16-84\%: -0.003 to +0.343); 83\% of the systematic MC shows a shortfall. Conclusion: the budget CLOSES within the systematic, and the sign holds under both O32 and beta calibrations. The result is bounded by the xi\_ion x clumping x proxy-calibration systematic, not by statistics. Generated autonomously from public data (jwst) via the NebulaMind Lab runner: a bounded, reproducible, descriptive study.
\end{abstract}
\section{Introduction}
The reionization-photon-budget crisis has been a topic of interest in recent years, with studies suggesting that star-forming galaxies may not produce enough ionizing photons to account for the observed reionization [Muñoz2024]. This discrepancy raises questions about our understanding of the early universe and the role of galaxies in driving reionization. Previous works have explored various aspects of this issue, including the calibration of excursion set reionization models [Park2022] and the assessment of galaxy ionizing photon budget at z < 10 [Duncan2015]. However, a comprehensive analysis of the reionization-photon-budget is still needed to reconcile these findings.
\section{Data and method}
To address this issue, we perform a literature-anchored budget calculation using the cosmic star formation rate density (SFRD) from Madau \& Dickinson (2014). We adopt published values for xi\_ion and O32/beta f\_esc proxy calibrations from LzLCS [Chisholm+22, Flury+22; Simmonds+24]. Our method focuses on the ionizing-photon-budget reconciliation at z\~{}8, considering factors such as clumping (C=2-5) and the JWST-SFRD tail.
\section{Result}
Our analysis reveals that star-forming galaxies require a escape fraction of f\_esc=0.210 (+0.211/-0.107) to close the reionization ionizing-photon-budget at z\~{}8, assuming the Madau-Dickinson SFRD, log xi\_ion=25.5+/-0.15, and clumping C=2-5. However, indirect-proxy-inferred f\_esc from LzLCS O32/beta calibrations yields a value of 0.062 (+0.108/-0.039). The median delta between the required and inferred values is +0.130 dex-frac (16-84\%: -0.003 to +0.343), with 83\% of systematic Monte Carlo simulations showing a shortfall.
\begin{figure}\centering\includegraphics[width=\columnwidth]{result.png}\caption{The reionization ionizing-photon-budget shortfall for star-forming galaxies is not robust to systematics at z\~{}8}\end{figure}
\section{Caveats}
It is essential to acknowledge the limitations of our approach, which relies on automated, single-selection, uncalibrated measurements. The accuracy of our results depends heavily on the assumptions made in the literature-anchored budget calculation and the adopted values for xi\_ion and O32/beta f\_esc proxy calibrations. Additionally, our analysis does not account for potential systematic errors or uncertainties in the underlying data, which may impact the validity of our conclusions. Further studies are needed to refine these estimates and provide a more comprehensive understanding of the reionization-photon-budget crisis.
\section*{References}
\footnotesize
[Muoz2024] Muñoz et al. (2024). Reionization after JWST: a photon budget crisis?. bibcode 2024MNRAS.535L..37M \par [Park2022] Park et al. (2022). Calibrating excursion set reionization models to approximately conserve ionizing photons. bibcode 2022MNRAS.517..192P \par [Duncan2015] Duncan et al. (2015). Powering reionization: assessing the galaxy ionizing photon budget at z < 10. bibcode 2015MNRAS.451.2030D \par [Davies2021] Davies et al. (2021). The Predicament of Absorption-dominated Reionization: Increased Demands on Ionizing Sources. bibcode 2021ApJ...918L..35D \par [Madau2017] Madau (2017). Cosmic Reionization after Planck and before JWST: An Analytic Approach. bibcode 2017ApJ...851...50M

\end{document}
